\documentclass[twoside]{article}

\usepackage{aistats2026}
\usepackage{amsmath,amssymb}
\usepackage{enumitem}
\usepackage{graphicx}
\usepackage{hyperref}

% If your paper is accepted, change the options for the package
% aistats2026 as follows:
%
%\usepackage[accepted]{aistats2026}
%
% This option will print headings for the title of your paper and
% headings for the authors names, plus a copyright note at the end of
% the first column of the first page.

% We also include a `preprint' option for non-anonymous preprints. 
% Change the options for the package aistats2026 as follows:
%
%\usepackage[preprint]{aistats2026}
%
% This option will print headings for the title of your paper and
% headings for the authors names, but does not print the copyright and 
% venue note at the end of the first column of the first page.

% If you set papersize explicitly, activate the following three lines:
%\special{papersize = 8.5in, 11in}
%\setlength{\pdfpageheight}{11in}
%\setlength{\pdfpagewidth}{8.5in}

% If you use the natbib package, activate the following three lines:
%\usepackage[round]{natbib}
%\renewcommand{\bibname}{References}
%\renewcommand{\bibsection}{\subsubsection*{\bibname}}

% If you use BibTeX in apalike style, activate the following line:
%\bibliographystyle{apalike}

\begin{document}

% If your paper is accepted and the title of your paper is very long,
% the style will print as headings an error message. Use the following
% command to supply a shorter title of your paper so that it can be
% used as headings.
%
%\runningtitle{I use this title instead because the last one was very long}

% If your paper is accepted and the number of authors is large, the
% style will print as headings an error message. Use the following
% command to supply a shorter version of the author names so that
% they can be used as headings (for example, use only the surnames)
%
%\runningauthor{Surname 1, Surname 2, Surname 3, ...., Surname n}

\twocolumn[

\aistatstitle{What Else Should Be True? Verifying Scientific Hypotheses from Indirect Evidence}

\aistatsauthor{ Author 1 \And Author 2 \And  Author 3 }

\aistatsaddress{ Institution 1 \And  Institution 2 \And Institution 3 } ]

\begin{abstract}
Hypothesis generation is becoming increasingly accessible to large language models and scientific agents, shifting the bottleneck in automated discovery towards deciding which candidate hypotheses are most valuable to pursue experimentally. This is particularly difficult due to the intrinsic tension between novelty and support: presence of strong direct evidence may imply that a hypothesis is already established, while the absence of such evidence does not immediately imply that it is false. We propose to address this tension by verifying hypotheses through their consequences. Our verifier asks: ``If (H) is true, what else should be true?'' Our verifier searches the scientific literature for evidence concerning these derived consequences, which may lie at different inferential distances from the original hypothesis. We organize these consequences in an explicit consequence graph and aggregate the retrieved evidence with a Bayesian interpretation. This yields a decomposable and auditable alternative to holistic LLM judging or debate, designed to recover support for novel hypotheses even when direct evidence is sparse.
\end{abstract}

\section{Introduction}

\paragraph{How should we verify scientific hypotheses?}
Recent advances in large language models and scientific agents have made it
increasingly easy to generate plausible scientific hypotheses. As hypothesis
generation scales, however, the bottleneck shifts from proposing candidate
ideas to deciding which of them are most likely to be correct and therefore
worth pursuing experimentally. A useful scientific discovery system should not
only generate hypotheses, but also distinguish promising candidates from
plausible but incorrect alternatives before expensive experiments are carried
out.

We study scientific hypothesis verification as a ranking problem. Given a scientific question and a set of competing hypotheses the goal is to rank the hypotheses according to how likely they are to be valid using only scientific knowledge available prior to their eventual direct validation or refutation. This setting reflects a realistic scientific decision problem: experimental resources are limited, and the value of a hypothesis generating system depends on whether promising hypotheses can be distinguished from plausible but incorrect alternatives before expensive experiments are performed.

\begin{figure*}[t]
    \centering
    \includegraphics[width=\linewidth]{figures/figure1_chatgpt.png}
    \caption{\textbf{Illustrative consequence graph for scientific hypothesis
    verification.} Three competing hypotheses about mutation $M$ imply
    overlapping mechanistic and observable consequences at different
    inferential distances. Literature evidence can attach to any empirically
    assessable proposition in the graph, providing direct or indirect evidence
    for the competing hypotheses. Colors denote hypothesis-associated pathways;
    purple nodes are shared consequences. Literature entries are illustrative
    placeholders, not real citations.}
    \label{fig:consequence-graph}
\end{figure*}

\paragraph{Novel hypotheses are difficult to verify directly.}
A central difficulty is that novelty and direct verifiability are naturally in
tension. If a hypothesis is genuinely novel, the existing literature is
unlikely to contain strong direct evidence establishing the same claim;
otherwise, the hypothesis would already be close to established knowledge. At
the same time, absence of direct evidence does not imply that the hypothesis is
unsupported. Scientific hypotheses make commitments about the world, and many
of their consequences may already have been studied independently and for
unrelated purposes. Evidence relevant to a novel hypothesis may therefore be
distributed across its consequences, even when no paper explicitly discusses
the hypothesis itself.

This observation creates a distinct verification problem from standard claim matching. Evidence may occur at different inferential distances from the focal hypothesis, and its relevance depends on how strongly the corresponding observation is implied by that hypothesis. Moreover, a consequence that is compatible with all competing hypotheses provides little evidence for ranking them, even if the consequence itself is well established. Verification must therefore identify evidence that is not only relevant, but discriminative among plausible alternatives. Figure~\ref{fig:consequence-graph} illustrates this setting for three competing hypotheses. Each hypothesis induces a branching set of consequences, with some pathways shared across candidates and literature evidence attaching to different parts of the graph.


\paragraph{Beyond holistic hypothesis evaluation.}
Existing approaches largely evaluate scientific hypotheses holistically. A
language model can be asked to directly score candidate hypotheses according to
plausibility or expected experimental success, while debate or tournament-style
procedures encourage explicit comparison between competing candidates. These
approaches can improve hypothesis ranking, but they place substantial burden on
the model's latent scientific knowledge and reasoning ability: the model must
implicitly identify which facts matter, determine how they bear on each
hypothesis, and combine them into a final judgment.

Retrieval-augmented approaches provide a complementary strategy by grounding evaluation in the scientific literature. When relevant evidence directly concerns the focal hypothesis, retrieving and assessing that evidence is natural. For a novel hypothesis, however, the informative literature may not discuss the hypothesis itself. As Figure~\ref{fig:consequence-graph} illustrates, relevant observations may instead be distributed across downstream propositions that are only indirectly connected to the focal claim. The challenge is therefore partly one of evidence formulation. Rather than asking only, ``Which papers discuss this hypothesis?'', we should ask, ``What else should hold if this hypothesis were correct, and has any of it already been observed?''

\paragraph{Verification through consequences.}
We propose to decompose hypothesis verification into a sequence of smaller,
explicitly grounded tasks. Given competing hypotheses, we construct scientific
consequences whose expected outcomes differ across candidates, search the
literature for evidence concerning these consequences, assess the strength and
direction of the retrieved evidence, and aggregate this evidence into a ranking
of the original hypotheses. Consequences may occur at different inferential
distances from the focal claim: direct evidence is used when available, while
more remote evidence can provide weaker but complementary support.

The resulting procedure has a natural Bayesian interpretation. If a
proposition $X$ is substantially more likely to hold under $H_i$ than under
$H_j$, then evidence supporting $X$ should increase the relative support for
$H_i$ over $H_j$. This perspective shifts the goal from accumulating facts that are merely
compatible with a hypothesis to identifying observations that change the
relative evidence for competing hypotheses. It also yields a decomposable
verification trace, in which the final ranking can be attributed to explicit
consequences, retrieved evidence, and inferential relationships instead of
being collapsed into a single holistic model judgment.

\paragraph{Contributions.}
$\blacktriangleright$~\textbf{Scientific hypothesis verification.}
We formulate the problem of ranking competing scientific hypotheses using only
scientific evidence available prior to their eventual direct validation.
$\blacktriangleright$~\textbf{Consequence-based evidence discovery.}
We introduce a verification framework that searches beyond direct mentions of
the focal hypotheses by identifying discriminative scientific consequences at
multiple inferential distances.
$\blacktriangleright$~\textbf{Structured evidence aggregation.}
We develop a structured representation of inferential relationships that
supports evidence accumulation while distinguishing evidence strength,
dependence, contradiction, and missing evidence.
$\blacktriangleright$~\textbf{Decomposable verification.}
Our framework exposes the intermediate consequences and retrieved evidence that
drive each ranking decision, allowing failures of consequence generation,
retrieval, evidence assessment, and aggregation to be analysed separately.
$\blacktriangleright$~\textbf{Empirical evaluation.}
We compare consequence-based verification with direct LLM judging,
debate-based ranking, direct retrieval-augmented verification, and unstructured
scientific search.

\begin{table*}[h]
\centering
\caption{\textbf{Comparison of methods for scientific hypothesis verification
and ranking.} $\checkmark$ denotes a capability explicitly represented by the
method, $\circ$ denotes partial or implicit support, and $\times$ denotes that
the capability is not explicitly modeled.}
\scriptsize
\setlength{\tabcolsep}{5pt}
\vspace{0.5em}
\begin{tabular}{lcccccc}
\hline
\textbf{Method}
& \shortstack{Literature\\grounding}
& \shortstack{Competing\\hypotheses}
& \shortstack{Indirect\\consequences}
& \shortstack{Structured\\evidence}
& \shortstack{Missing\\$\mathbf{\neq}$refuting}
& \shortstack{Decomposable\\trace} \\
\hline
Prompted LLM ranking \cite{liu2026researchbench}
    & $\times$ & $\checkmark$ & $\times$ & $\times$ & $\times$ & $\times$ \\
Logit-based energy scoring \cite{rajwal2026logit}
    & $\times$ & $\checkmark$ & $\times$ & $\times$ & $\times$ & $\times$ \\
Comparative empirical forecasting \cite{mule2026forecasting}
    & $\times$ & $\checkmark$ & $\times$ & $\times$ & $\times$ & $\circ$ \\
Co-Scientist Ranking \cite{gottweis2026coscientist}
    & $\times$ & $\checkmark$ & $\times$ & $\times$ & $\times$ & $\circ$ \\
Co-Scientist Deep Verification \cite{gottweis2026coscientist}
    & $\checkmark$ & $\times$ & $\circ$ & $\circ$ & $\times$ & $\checkmark$ \\
\textbf{Consequence-based verification (ours)}
    & $\checkmark$ & $\checkmark$ & $\checkmark$ & $\checkmark$
    & $\checkmark$ & $\checkmark$ \\
\hline
\end{tabular}
\label{tab:related-work}
\end{table*}

\section{Problem Setup}
\label{sec:problem}

Let $Q$ denote a scientific question and let
$\mathcal{H}=\{H_1,\ldots,H_k\}$ denote a set of competing candidate
hypotheses. We assume that the hypotheses are sufficiently novel that novelty
itself is not part of the ranking objective. Let $\mathcal{L}_{\leq t}$ denote
the scientific literature available up to a cutoff time $t$, chosen prior to
the direct experimental validation or refutation used to evaluate the
hypotheses.

A verifier receives $(Q,\mathcal{H},\mathcal{L}_{\leq t})$ and assigns each
$H_i\in\mathcal{H}$ a scalar score $S(H_i)$, inducing a ranking over the
candidate set. The score is intended to reflect the evidence for $H_i$ in the
pre-cutoff scientific record: hypotheses receiving higher scores should be more
likely to be subsequently validated. The temporal restriction is essential to
the task, since evidence appearing after $t$ may directly reveal the eventual
outcome and is reserved for evaluation.

\paragraph{Evidence through scientific consequences.}
The hypotheses need not be directly discussed in $\mathcal{L}_{\leq t}$.
Instead, a hypothesis can bear on a collection of scientific propositions at
different inferential distances. For a proposition node $v$, let $X_v$ denote
the binary variable indicating whether that proposition holds. Under competing
hypotheses, the same proposition may have different predictive probabilities,
\begin{equation}
    P(X_v=1\mid H_1),\ldots,P(X_v=1\mid H_k).
\end{equation}
Evidence concerning $X_v$ is discriminative whenever it changes the relative
support assigned to the candidate hypotheses. In particular, an observation
that is equally likely under all hypotheses carries little ranking information,
even if it is strongly established in the literature.

\paragraph{The consequence graph.}
We represent these relationships using a consequence graph $G=(V,E)$, as
illustrated in Figure~\ref{fig:consequence-graph}. Nodes represent scientific
propositions and directed edges represent inferential relationships between
them. Candidate hypotheses form distinguished root nodes, while empirically
searchable propositions may occur at any depth in the graph. Different
hypotheses may imply shared consequences, and evidence may therefore bear on
several candidates simultaneously. Graph depth is not itself a measure of
evidence strength: a distant consequence can remain highly informative when
the intervening implications are strong, while a nearby consequence can be only
weakly diagnostic.

Let $D$ denote the scientific evidence retrieved for propositions in the graph.
At the level of two competing hypotheses, evidence changes their posterior odds
according to
\begin{equation}
\frac{P(H_a\mid D)}{P(H_b\mid D)}
=
\frac{P(H_a)}{P(H_b)}
\frac{P(D\mid H_a)}{P(D\mid H_b)}.
\end{equation}
Our objective is to construct and verify propositions for which this likelihood
ratio is informative, and to combine the resulting evidence into scores
$S(H_i)$ that accurately rank the candidates.

\paragraph{Desiderata.}
An ideal verifier should satisfy the following properties.
\begin{itemize}[left=0.1cm]
    \item \textbf{Exploit evidence at multiple inferential distances.}
    Direct evidence should be used when available, but the verifier must remain
    effective when the focal hypothesis has not yet been explicitly tested.
    Relevant evidence may concern consequences several inferential steps removed
    from the hypothesis; its contribution should depend on the strength of the
    corresponding implications rather than graph distance alone.

    \item \textbf{Prioritise discriminative evidence.}
    Evidence is useful for ranking hypotheses only to the extent that it changes
    their relative support. A proposition that is equally compatible with all
    candidates provides little information, even if it is strongly established.

    \item \textbf{Avoid search-policy-induced confirmation bias.}
    Evidence acquisition should be hypothesis-symmetric and should not depend on
    whether previously retrieved evidence happened to favour a particular
    candidate. Search queries should target empirical propositions or
    observations without encoding a desired conclusion, and retrieved outcomes
    should not be selectively retained according to whether they support a
    preferred hypothesis.

    \item \textbf{Account for evidence strength and dependence.}
    Multiple strong and independent studies should contribute more than a single
    weak observation, while correlated measurements or papers reflecting the
    same underlying mechanism should not be treated as fully independent
    confirmations.

    \item \textbf{Distinguish contradiction from missing evidence.}
    Many relevant consequences of a novel hypothesis may simply not have been
    studied. Failure to retrieve informative evidence should therefore not be
    interpreted as evidence against a hypothesis.

    \item \textbf{Support decomposable and auditable verification.}
    The final ranking should be traceable to explicit consequences, retrieved
    evidence, and inferential relationships rather than arising from a single
    holistic model judgment. This allows errors in consequence construction,
    retrieval, evidence assessment, and aggregation to be identified
    separately.
\end{itemize}

Together, these desiderata define the focus of this work: to recover reliable
ranking signal from scientific evidence that may be distributed across direct
and indirect consequences, while retaining an explicit account of how that
evidence bears on the competing hypotheses.

\section{Related Work}
\label{sec:related-work}

\paragraph{Scientific hypothesis ranking and forecasting.}
Recent work has begun to study the evaluation of scientific hypotheses as a
distinct problem from hypothesis generation. ResearchBench
\cite{liu2026researchbench} formulates hypothesis ranking as one of three core
subtasks of scientific discovery and evaluates whether LLMs can distinguish a
gold hypothesis from plausible alternatives. Subsequent work has explored
alternatives to prompted LLM judging. Rajwal et al.
\cite{rajwal2026logit} score candidate hypotheses using intrinsic
logit-based model confidence, while Mule et al.
\cite{mule2026forecasting} train language models to predict which of two
research ideas will achieve stronger empirical performance. These approaches
treat verification primarily as a direct prediction problem: the final ranking
is produced from the candidate hypotheses and the model's internal
representation or learned judgment. Our setting instead asks whether the
pre-existing scientific record contains evidence that should change the
relative support for competing hypotheses.

\paragraph{Agentic hypothesis evaluation.}
Co-Scientist \cite{gottweis2026coscientist} provides the closest agentic
approach to our setting. Its Ranking agent compares hypotheses through
pairwise scientific debates and aggregates the resulting preferences using an
Elo tournament. Its Reflection agent additionally performs literature-grounded
reviews, including a deep verification procedure that decomposes a hypothesis
into assumptions and subassumptions, and an observation review that asks
whether the hypothesis can explain existing experimental findings. Our method
shares the motivation of decomposing verification into smaller scientific
judgments, but differs in how evidence is represented and combined. We
construct a shared graph of propositions predicted by the competing
hypotheses, explicitly assess the strength of both inferential relationships and
retrieved evidence, and propagate these quantities back to a posterior score.
This makes indirect evidence, shared consequences, evidence dependence, and
missing evidence explicit parts of the verification model.

\begin{figure*}[t]
    \centering
    \includegraphics[width=\textwidth]{figures/figure2_method_pipeline.png}
    \caption{\textbf{Overview of the proposed verification pipeline.}
    Starting from a scientific question, competing hypotheses, and pre-cutoff
    literature, the verifier constructs a shared consequence graph, assigns
    ordinal strengths to its inferential relationships, retrieves and assesses
    evidence for empirically assessable propositions, and propagates the
    calibrated evidence through the graph to rank the candidate hypotheses.}
    \label{fig:method-overview}
\end{figure*}

\paragraph{Literature-grounded verification and indirect discovery.}
A broader line of work studies whether scientific claims can be verified
against retrieved literature. Recent evaluations of medical causal hypothesis
verification, for example, show that LLMs can struggle to provide valid
supporting articles and to reject unsupported hypotheses
\cite{ahmed2026medical}. Literature-based discovery provides an earlier
perspective on indirect evidence, showing that scientifically useful
relationships can emerge by connecting findings distributed across otherwise
separate literatures \cite{swanson1986undiscovered}. These works motivate
literature grounding and multi-step scientific reasoning, but address a
different objective. Claim verification asks whether evidence supports a focal
claim, while literature-based discovery searches for previously unrecognized
connections. We instead use distributed literature evidence to discriminate
among a fixed set of competing scientific hypotheses.


\section{Method}
\label{sec:method}

\subsection{Overview}
Given a scientific question $Q$, a set of competing hypotheses
$\mathcal{H}={H_1,\ldots,H_k}$, and access to the pre-cutoff literature
$\mathcal{L}_{\leq t}$, our goal is to rank the hypotheses using evidence that
may be distributed across their scientific consequences. Our verifier consists
of four stages. First, we expand the hypotheses into a shared
\emph{consequence graph} containing propositions at different inferential
distances from the candidate hypotheses. Second, an LLM assigns ordinal
strengths to the inferential relationships in the graph. Third, for each
empirically assessable proposition, we retrieve the relevant literature. An evidence assessor summarizes the
direction and strength of the evidence with respect to the proposition. Finally, we
map the ordinal judgments to probabilistic quantities and propagate the
evidence through the graph to obtain a posterior score for each hypothesis.
Figure~\ref{fig:method-overview} summarizes the complete pipeline.

We use a single categorical random variable
$H\in{1,\ldots,k}$ to represent which candidate hypothesis is correct, with a
uniform prior unless stated otherwise. Each non-hypothesis node $v$ in the
consequence graph is associated with a binary latent variable $X_v$ indicating
whether the corresponding scientific proposition holds. The retrieved
literature evidence $D_v$ provides a noisy observation of $X_v$ for those
propositions for which informative evidence is available. The resulting
decomposition separates two questions that are difficult to answer reliably in
a single holistic judgment:
\begin{equation} 
\underbrace{P(X_v=1\mid H)}_{\text{what should follow if $H$ is true?}}, \quad \text{and} 
\end{equation} 
\begin{equation} 
\underbrace{P(D_v\mid X_v)}_{\text{what does the literature say about $X_v$?}}. \end{equation}

\subsection{Constructing the Consequence Graph}
\label{sec:graph-construction}

\paragraph{Recursive consequence generation.}
We construct a directed graph $G=(V,E)$ of consequences. Starting from each candidate hypothesis $H_i$, we recursively derive scientific propositions that should become more or less likely if the parent proposition holds:
\begin{equation}
H_i \rightarrow X_1 \rightarrow \cdots \rightarrow X_n.    
\end{equation}
The resulting nodes need not be separated into ``intermediate'' and ``observable'' propositions: any node in the graph may be empirically assessable and therefore associated with retrieved literature evidence.

\paragraph{Shared and discriminative consequences.}
The graph is shared across hypotheses; we do not construct $k$
independent trees. We pool propositions generated from different hypotheses
and merge semantically equivalent nodes, allowing inferential pathways to
collide as in Figure~\ref{fig:consequence-graph}. Such collisions are important:
a proposition predicted by several hypotheses is less discriminative than one
whose expected truth differs across candidates, while a shared intermediate
node also makes explicit when several downstream observations reflect the same
underlying mechanism. During generation, the model is therefore encouraged to
propose consequences that distinguish the focal hypothesis from at least one
alternative, while also exploring distinct inferential pathways and avoiding
near-duplicate manifestations of a single mechanism.

\paragraph{Outcome-blind construction.}
To avoid selecting search directions based on favorable evidence, graph
construction is completed before literature outcomes are observed. The
generation and merging procedure depends on $(Q,\mathcal{H})$ and the model's
scientific reasoning, but not on whether later retrieval supports or
contradicts any candidate. Once constructed, the graph is frozen and literature
is queried for all retained empirically assessable nodes. Exact generation
budgets, stopping criteria, and semantic merging details are deferred to the
appendix.

\subsection{Parameterizing Inferential Relationships}
\label{sec:implication}

The edges of $G$ encode uncertain scientific implications, not
deterministic logical rules. For each edge $(u,v)\in E$, an LLM judges how
strongly the truth of the parent proposition implies the child proposition.
We ask the model for an ordinal judgment instead of a precise probability: 
$r_{uv}\in \{\textsc{strongly implied},$
$\textsc{implied},$
$\textsc{weakly implied},$
$\textsc{neutral},$
$\textsc{unlikely},$
$\textsc{strongly contradicted}\}$.
    
These labels define a placeholder ordered implication scale; the exact
categories are specified in the appendix. Ordinal judgments are easier to
interpret and allow the numerical scale to be calibrated separately from the
scientific reasoning that produced the graph.

A monotone mapping $f_{\theta}$ converts these labels into local probabilistic
parameters used by the graphical model. Conceptually,
\begin{equation}
    f_{\theta}(r_{uv})
    \approx
    P(X_v=1\mid X_u=1),
\end{equation}
with stronger implication labels mapped to larger probabilities and
contradictory labels mapped below the neutral level. When a proposition has
multiple parents, the corresponding local conditional distribution combines
the incoming calibrated relations using a fixed monotone composition rule. We
describe the exact parameterization in the appendix; the main requirement is
that stronger ordinal implications induce stronger probabilistic dependence
and that shared descendants are represented only once in the graph.

\subsection{Retrieving and Assessing Literature Evidence}
\label{sec:evidence}

\paragraph{Outcome-neutral retrieval.}
Any node whose proposition is empirically assessable can be queried, regardless
of its distance from the hypotheses. For each such node $v$, we generate
literature queries that describe the underlying entities, intervention or
comparison, and measured quantity without encoding the outcome favored by a
particular hypothesis. For example, the query targets studies measuring a marker in $M$-mutant and
comparison systems without asserting in advance that the marker
\emph{increases}. Retrieval is restricted to
$\mathcal{L}_{\leq t}$.

\paragraph{Ordinal evidence assessment.}
Let $D_v$ denote the set of retrieved studies relevant to proposition $X_v$.
An evidence assessor reads the retrieved evidence jointly and returns an
ordinal judgment
$s_v\in \{\textsc{strong support},$
$\textsc{support},$
$\textsc{weak support},$
$\textsc{mixed},$
$\textsc{weak contradiction},$
$\textsc{contradiction},$
$\textsc{strong contradiction},$
$\textsc{no evidence}\}$.

These labels define a placeholder evidence scale. We keep
\textsc{mixed} and \textsc{no evidence} as distinct outcomes, even if both may
receive little weight after calibration. The assessment is instructed to
account for the relevance and directness of the studies, their
methodological strength, agreement across results, and effective independence.
This allows several strong independent studies to provide more evidence than a
single weak observation, while avoiding a simple paper-counting rule for
correlated reports of the same underlying result.
We map $s_v$ to an evidence log-likelihood ratio
\begin{equation}
    g_{\phi}(s_v)
    =
    \log
    \frac{P(D_v\mid X_v=1)}
         {P(D_v\mid X_v=0)}.
    \label{eq:evidence-lr}
\end{equation}
Strong support yields a positive value, contradiction a negative value, and
mixed or uninformative evidence lies near zero. In particular, when no
informative literature is found, the corresponding likelihood ratio is set
close to one, so that missing evidence does not by itself count against a
hypothesis.

\subsection{Bayesian Evidence Propagation}
\label{sec:bayesian-propagation}

The consequence graph together with the calibrated implication and evidence
parameters defines a probabilistic model over the candidate hypothesis and
latent scientific propositions. Let $\mathbf{X}$ denote all proposition
variables and let $\mathcal{O}\subseteq V$ denote the nodes for which literature
evidence is available. The joint distribution factorizes as
\begin{equation}
\begin{split}
    P_{\theta,\phi}(H,\mathbf{X},D\mid G)
    &=
    P(H)
    \prod_{v\in V}
    P_{\theta}(X_v\mid X_{\mathrm{pa}(v)},H) \\
    &\qquad\cdot
    \prod_{v\in\mathcal{O}}
    P_{\phi}(D_v\mid X_v),
\end{split}
\label{eq:joint-model}
\end{equation}
where the first product encodes the inferential structure of the consequence
graph and the second attaches literature evidence to observed propositions.

The score assigned to hypothesis $H_i$ is its posterior probability after
marginalizing over the unobserved propositions,
\begin{equation}
\begin{split}
    S(H_i)
    &=
    P_{\theta,\phi}(H=i\mid D,G)  \\
    &\propto
    P(H=i)
    \sum_{\mathbf{X}}
    P_{\theta,\phi}(\mathbf{X},D\mid H=i,G).
\end{split}
    \label{eq:hypothesis-posterior}
\end{equation}
The hypotheses are ranked by $S(H_i)$. This marginalization is what allows
evidence at different inferential distances to contribute coherently. Evidence
for a remote proposition is attenuated by uncertainty in the intervening
relations, while several observations descending from the same intermediate
mechanism update that shared latent proposition, so they are not treated as
independent confirmations of the root hypothesis. Unobserved nodes are simply
marginalized out.

\subsection{Validation-Set Calibration}
\label{sec:calibration}

The probabilistic interpretation requires numerical values for the ordinal
implication and evidence categories, but we do not train a separate verifier.
Instead, we initialize the monotone mappings $f_{\theta}$ and $g_{\phi}$ with
fixed values and calibrate only this small set of global parameters on a held-out
validation set. All LLM prompts, generated graphs, retrieval results, and
ordinal judgments are frozen before calibration.

Let $\mathcal{V}$ denote the validation set and let $H_n^\star$ be the
subsequently validated hypothesis for instance $n$. We fit the calibration
parameters by minimizing the negative log posterior of the gold hypothesis,
\begin{equation}
    (\theta^\star,\phi^\star)
    =
    \arg\min_{\theta,\phi}
    \sum_{n\in\mathcal{V}}
    -\log
    P_{\theta,\phi}
    \left(
        H_n^\star
        \mid
        D_n,G_n
    \right),
    \label{eq:calibration-objective}
\end{equation}
subject to monotonicity constraints that preserve the ordering of the ordinal
categories. Calibration therefore adjusts how much probabilistic weight is
assigned to judgments such as ``strong implication'' or ``moderate support''
without changing which consequences are generated or which evidence is
retrieved. We evaluate both the fixed zero-shot mappings and the
validation-calibrated variant in our experiments.

\section{Experiments}

\section{Discussion}

\clearpage

\bibliography{references}
\bibliographystyle{abbrv}

\end{document}
